\begin{table}

\caption{\label{tab:h1-capacidade}Rendimento potencial municipal no cenário de alto insumo e sequeiro}
\centering
\begin{tabular}[t]{l|l|r|r|r|l}
\hline
Cultura & Predominancia & Municipios & N válido & Mediana & P25--P75\\
\hline
Milho & Amazônia & 76 & 76 & 3002 & 2.862--3.264\\
\hline
Milho & Cerrado & 61 & 61 & 3320 & 2.797--3.572\\
\hline
Milho & Pantanal & 4 & 4 & 6790 & 6.321--6.933\\
\hline
Soja & Amazônia & 76 & 76 & 2701 & 2.557--2.852\\
\hline
Soja & Cerrado & 61 & 61 & 3261 & 2.817--3.390\\
\hline
Soja & Pantanal & 4 & 4 & 3421 & 3.386--3.457\\
\hline
\end{tabular}
\end{table}
\par\smallskip{\footnotesize Fonte: elaboração própria com GAEZ v4, conforme o dicionário fornecido. Unidade: kg de matéria seca por hectare; cada município tem peso igual. N válido informa a cobertura de cada resumo. O cenário não é rendimento realizado e não garante capacidade predeterminada em relação a todo o painel.\par}

